\documentclass[
  paper=a4,
  fontsize=12pt,
  DIV=16,
  headheight=52pt,
  footheight=45pt,
  headinclude,
  parskip=full,
]{scrartcl}

\usepackage{fontspec}
\setsansfont{Source Sans Pro}
\renewcommand{\familydefault}{\sfdefault}

\usepackage{polyglossia}
\setmainlanguage{german}

\usepackage{microtype}

\usepackage{scrlayer-scrpage}
\pagestyle{scrheadings}
\setkomafont{pagehead}{\normalfont}
\setkomafont{pagefoot}{\normalfont\footnotesize}
\ohead{\includegraphics[height=1.5cm]{pep.pdf}}
\ihead{%
  \large\bfseries PeP et al.\ e.\,V.\\%
  Begabtenförderungsprogramm –\\
  Initiativstipendien%
}
\cfoot{}
\ifoot{%
  PeP et al.~e.\,V.\\
  Begabtenförderung\\
  \url{www.pep-dortmund.org}
}


\usepackage{titling}
\usepackage{booktabs}

\usepackage{xcolor}
\usepackage{calc}
\usepackage[colorlinks=true,urlcolor=blue!50!black]{hyperref}

\newcommand\MyTextField[2][]{\TextField[#1, backgroundcolor=black!10, charsize=0pt, borderwidth=0]{#2}}
\renewcommand*{\LayoutTextField}[2]{\makebox[\widthof{#1: }][l]{#1: }%
\raisebox{0.8\baselineskip}{\raisebox{-\height}{#2}}}


\begin{document}
\section*{Bewerbungsunterlagen}
\emph{Formale Hinweise}

Bitte reiche Deine Bewerbung bis zum Quartalsende (31. März, 30. Juni,
30. September oder 31. Dezember) beim PeP-Vorstand
(Henning Moldenhauer, Alex Birnkraut oder Vanessa Müller) ein,
addressiert an:

PeP et al. e.V.\\
Henning Moldenhauer\\
CP-01-190\\[0.5\baselineskip]
Technische Universität Dortmund\\
Otto-Hahn-Str.~4a\\
44221 Dortmund

Um für das Bewerbungsverfahren zugelassen zu werden, muss Dein Antrag bestehen aus:
\begin{itemize}
  \item einem vollständigen, tabellarischen \textbf{Lebenslauf},
  \item einem 1- bis 2-seitigen \textbf{Motivationsschreiben}, in welchem Du
	  begründest warum Dein Projekt förderungswürdig ist,
  \item eine \textbf{Selbstdarstellung}, aus welcher dein soziales/
	  gesellschaftliches/fakultätsinternes und/oder sonstiges Engagement hervorgeht.
  \item einer \textbf{Aufschlüsselung} Deiner kosten, sodass wir abschätzen können,
	  ob eine Förderung notwendig ist.
\end{itemize}

\textbf{%
  Eine Bewerbung ist nur möglich, wenn du an der Technischen Universität Dortmund
  als Student eingeschrieben bist und Physik im Studiengang
  \emph{Bachelor of Science Physik},
  \emph{Bachelor of Science Lehramt Physik (BfP)},
  \emph{Bachelor of Science Medizinphysik},
  \emph{Master of Science Physik},
  \emph{Master of Science Lehramt Physik (MfP)} oder
  \emph{Master of Science Medizinphysik} studierst.
}

Bitte reiche nur DIN A4 Dokumente ein. Klammer oder hefte diese nicht und nutze auch
keine Klarsichtfolien. Unvollständige Bewerbungen werden nicht berücksichtigt.

Für persönliche Fragen stehen Dir Thorben Menne (CP-03-142) und
Henning Moldenhauer (CP-O1-190) zur Verfügung.
\end{document}
